\documentclass{beamer}
\usepackage[utf8]{inputenc}
\usepackage[T1]{fontenc}
\usepackage[portuguese]{babel}
\usepackage{amsmath}
\usepackage{tikz}
\usetikzlibrary{shapes.geometric, arrows.meta, 3d}

% Configuração do Tema
\usetheme{Madrid}
\usecolortheme{default}

% Informações do Título
\title{Fundamentos da Disciplina}
\subtitle{Fundamentos de Fisica}

\date{}

\begin{document}
	
	% --- Slide 1: Título ---
	\begin{frame}
		\titlepage
	\end{frame}
	
	% --- Slide 2: Avaliação ---
	\begin{frame}{Método de Avaliação}
		\begin{block}{Componentes da Avaliação}
			\begin{table}[]
				\begin{tabular}{l l c}
					\textbf{Componente} & \textbf{Data} & \textbf{Peso} \\
					\hline
					A) Teste 1 & 13/04/26 & 40\% \\
					B) Teste 2 & 08/06/26 & 40\% \\
					C) Projeto/Trabalho & 15/06 e 22/06 & 20\% \\
				\end{tabular}
			\end{table}
		\end{block}
		
		\begin{alertblock}{Exame Oral Facultativo (D)}
			Para melhorar a nota final: de 0 a 6 pontos adicionais. É possível aceder ao exame oral apenas com nota mínima de 7.
		\end{alertblock}
		
		\vspace{0.5cm}
		\textbf{** Alternativa:}
		A) e B) é a frequência final que vale o 80\% da nota final no exame oral.
	\end{frame}
	
	% --- Slide 3: Cálculo da Nota Final ---
	\begin{frame}{Cálculo da Nota Final}
		\textbf{Fórmula:}
		$$ Nota_{Final} = (Nota_{Teste1} \cdot 0,4) + (Nota_{Teste2} \cdot 0,4) + (Proj \cdot 0,2) $$
		
		$$ Nota_{FinalTotal} = A + B + C + \text{pontos extra de } D $$
		
		\begin{exampleblock}{Exemplo de Cálculo}
			Supondo as notas: Teste 1 = 8; Teste 2 = 12; Projeto = 16.
			$$ (8 \cdot 0,4) + (12 \cdot 0,4) + (16 \cdot 0,2) $$
			$$ = 3,2 + 4,8 + 3,2 = 11,2 $$
			
			\textbf{Com Oral:} O estudante decide fazer exame oral e tem nota 3.
			$$ Nota_{Final} = 11,2 + 3 = 14,2 $$
		\end{exampleblock}
		
		\footnotetext{Email para dúvidas: p8346@ulusofona.pt}
	\end{frame}
	

	
\end{document}